\scanpage{76}
\mspage{69}
% archive sheet number: 75
en celle $R$ du second.

Je dis que l'on a de \unsure{ttes|toutes} façons
\[
R\leq\alpha(R'),
\]
\uline{sans} hypothèses d'admissibilité sur $u$, i.e. que la bijection inverse de $(\alpha\alpha)$ donne par passage aux quotients un hom.
\[
\operatorname{Hom}_W(u^*F,\unsure{G'|C'})\longleftarrow\operatorname{Hom}_W(F,u_*(G')),
\tag*{($\alpha$ bis)}
\]
forcément surjectif, et \inserted{de plus} fonctoriel en $F$, \deleted{$\unsure{G'}$} variables dans $\mathrm{HO}_W(I)$ et $\mathrm{HO}_W(I')$ respectivement). En effet, on définit la flèche $(\alpha\ \text{bis})$ comme composée
\diagram{
\begin{tikzcd}[ampersand replacement=\&,column sep=large,row sep=9em]
\operatorname{Hom}_W(F,u_*G')\arrow[r,"{\begin{gathered}\text{car }u^*\text{ est}\\\text{comp. avec }W(I),\\W(I')\text{ sur }M(I)\\\deleted{\text{tout}}\text{ entier, et}\\\text{on interprète les}\\\text{deux mbres comme des}\\\operatorname{Hom}\text{ dans }M(I)W(I)^{-1}\\\text{et }M(I')W(I')^{-1}\end{gathered}}"']\&\operatorname{Hom}_W(u^*F,u^*u_*G')\arrow[dr,"{\begin{gathered}\text{composition}\\\text{avec l'hom}\\\text{d'adjonction}\\u^*u_*G\to G\end{gathered}}"]\&\\
\&\&\operatorname{Hom}_W(u^*F,G')
\end{tikzcd}
}
Cela \deleted{signifie} implique donc que le foncteur \inserted{sur $\mathrm{HO}_W(I)$}
\[
F\longmapsto\operatorname{Hom}_{\mathrm{HO}_W(I')}(u^*F,G')
\]
est, sinon représentable, du moins quotient terme à terme du foncteur représenté par $u_*G'$. $^*$

% margin note written slanted
\margin{$^*$ Mais il n'est pas clair, sans hyp. suppl. sur la situation, que ce foncteur représentable puisse être choisi en dépendance fonctorielle de $G'$, envisagé dans $\mathrm{HO}_W(I')$.}
